\scanpage{67}
% archive sheet number: 66
part les cycles de longueur $1$ et ceux de longueur \unsure{2}
\[
\begin{aligned}
X_\sigma&=\prod_{i\in I_1(\sigma)}X_{ii}
\prod_{\deleted{\unsure{i'\in}}i\in I_2(\sigma)} X_{\{i,\sigma i\}}
\prod_{i\in I_*(\sigma)}X_{\{i,\sigma i\}}\\
\varepsilon_\sigma&=(-1)^{\unsure{\nu|v}_2+\deleted{\unsure{\nu_1}}\,\nu_*^+}
\end{aligned}
\]
où
\[
\left|\begin{aligned}
I_1(\sigma)&=I^\sigma\\
I_2(\sigma)&=\bigcup\,\text{cycles de card. }2\\
I_*(\sigma)&=\bigcup\,\text{cycles de card }\geq3\\
\nu_1&=\unsure{\text{nb}}\ \text{de cycles de card }1=\operatorname{card}\overbrace{I^\sigma}^{\inserted{I_1(\sigma)}}\deleted{\unsure{\mathrm{car}|\mathrm{card}}}\\
\nu_2&=\text{nb de cycles de card }2\\
\nu_*^+&=\text{nb de cycles de cardinal pair } >\unsure{2|4}\\
\nu_*&=\text{nb de }\rule{7em}{0.35pt}>2
\end{aligned}\right.
\]
on trouve donc
\[
\Delta=\sum_{\sigma\in\unsure{\mathfrak S|\mathscr S}_I}(-1)^{\nu_2+\nu_*^+}2^{\nu_1}
\underbrace{\prod_{i\in I_1(\sigma)}\deleted{X}Y_{\{i\}}\!
\left(\prod_{\substack{\deleted{\unsure{i\in I_2(\sigma)/\sigma}}\\H\in\dot I_2(\sigma)}}\deleted{X}Y_H\right)^{2}
\prod_{H\in\dot I_*(\sigma)}\deleted{X}Y_{H,\sigma}}_{\deleted{\unsure{X_\sigma}}\,Y_\sigma}
\]
où
\[
\deleted{X}Y_{H,\sigma}=\prod_{a\in A(H,\sigma)}X_a
\]
$A(H,\sigma)$ est l'ens des \uline{arêtes} du polygône combinatoire (orienté) dont l'ens des sommets est $H$, avec la structure polyg. orientée déduite de $\sigma|H$ \enspace $(\dot I_2(\sigma)=I_2(\sigma)/\sigma,\ \dot I_*(\sigma)=I_*(\sigma)/\sigma)$.

\deleted{Les monômes} \inserted{On peut} dire que \inserted{la donnée de} $\sigma$ \deleted{définit} \inserted{équivaut à celle d'} une \deleted{\unsure{str}} décomposition de $I$ en $I_1(\sigma)\amalg I_2(\sigma)\amalg I_3(\sigma)$ où $I_2(\sigma)$ est munie d'une involution ss pt fixe $\sigma_2$ (i.e. d'une \deleted{\unsure{de}} \unsure{rel.|rel|relation} d'équivalence à classes de cardinal $2$), \deleted{et} et $I_3(\sigma)$ d'une structure de
